\begin{figure}[b!]
\centering
\begin{tcolorbox}[
    colback=myboxcolor, %
    colframe=blue!40!black, %
    boxrule=1pt, %
    arc=4pt, %
    width=0.85\textwidth,
    title={Content overview},
    fonttitle=\bfseries
]
{\small
\textbf{Formal power series:} limits, division, derivatives, substitution, exponentials and logarithms, non-integer powers,  generalized binomial theorem, Laurent series, multivariate series, infinite products, multipliability, distributive laws.

\textbf{Permutations:} sign, cycle decomposition, inversions and Lehmer codes, length generating function.

\textbf{Determinants:} Cauchy--Binet formula, Desnanot--Jacobi identity, Dodgson condensation, Jacobi's complementary minor theorem, Vandermonde determinant.

\textbf{Alternating sums and signed counting:} alternating-sum cancellations, sign-reversing involutions, inclusion--exclusion, weighted inclusion--exclusion, Boolean Möbius inversion, derangements, surjection counting.

\textbf{Partitions and q-series:} partitions and Ferrers diagrams, conjugation, partition generating function, odd parts and distinct parts correspondence, Euler pentagonal number theorem, Jacobi triple product identity, divisor-sum recurrences.

\textbf{q-binomial coefficients:} q-integers, q-factorials, Gaussian binomial coefficients, q-binomial theorems, subspace-counting interpretation, classical limits.

\textbf{Symmetric functions:} symmetric polynomials, elementary complete homogeneous and power-sum families, Newton--Girard relations, monomial symmetric polynomials and basis theorem, Schur polynomials, skew Schur polynomials, alternants, Bender--Knuth involutions, Pieri rules, Jacobi--Trudi identities, Littlewood--Richardson rule.

\textbf{Lattice paths:} lattice path enumeration, Lindström--Gessel--Viennot lemma, weighted path model, nonintersecting path families.

\textbf{Applications:} domino tilings (height-2 Fibonacci recurrence, height-3 classification).
}
\end{tcolorbox}
\caption{Overview of the contents from the source textbook \citep{grinberg2025introduction}.}
\label{fig:contents}
\end{figure}
